% META: {"id":"1SPE-DERGLOBAL-EX-015","chapitre":"1SPE-DERIVATION-GLOBAL","type_objet":"exercice","capacites":["1SPE-DERIVATION-GLOBAL-C2"],"capacites_codes":["C2"],"methodes":["M2"],"parcours":2,"competences":["calculer"],"duree_min":12,"mode_creation":"ex_nihilo","sources_inspiration":[],"fichier_tex":"chapitres/1SPE-DERIVATION-GLOBAL/exercices/1SPE-DERGLOBAL-EX-015.tex","corrige_tex":"chapitres/1SPE-DERIVATION-GLOBAL/corriges/1SPE-DERGLOBAL-CO-015.tex","status":"approved","origin":"generated","status_history":["generated","audited","accepted","approved"],"release_acceptance":"RELEASE_OWNER_BATCH_ACCEPTANCE","acceptance_closure_digest":"sha256:9b3ccf9a81c5520fbb7e03b7b2d3e7bf2057a02834b6d908bf3be5f49f3b553f","acceptance_receipt_digest":"sha256:4fad86f0282e0fa4e0419cca1ad6379ae7af5a280c04a4a90880f90a1c044242"}
% BEGIN-VERIFY
% from sympy import symbols, diff, expand, factor, solve
% x = symbols('x')
% f = (x**2 - 3)*(x**2 + 1)
% fp = diff(f, x)
% assert expand(fp) == 4*x**3 - 4*x
% assert factor(fp) == 4*x*(x - 1)*(x + 1)
% sols = solve(fp, x)
% assert set(sols) == {-1, 0, 1}
% END-VERIFY
\begin{exercice}{1SPE-DERGLOBAL-EX-015}{2}{12}
Soit $f(x) = (x^2-3)(x^2+1)$ définie sur $\mathbb{R}$.
\begin{enumerate}
  \item Calculer $f'(x)$ en appliquant la règle du produit, puis factoriser le résultat.
  \item Résoudre $f'(x) = 0$.
  \item Dresser le tableau de variations de $f$.
\end{enumerate}
\end{exercice}
